\section*{Глава 1. Анализ предметной области}
\addcontentsline{toc}{section}{Глава 1. Анализ предметной области}
\stepcounter{section}

\subsection{Обзор существующих CI/CD-систем и инструментов автоматизации}

Современная разработка программного обеспечения все чаще строится вокруг практик непрерывной интеграции и непрерывной доставки (CI/CD). CI/CD-пайплайн представляет собой формализованную последовательность этапов, в которой исходный код автоматически проверяется, собирается, тестируется и подготавливается к развертыванию. Такой подход снижает количество ручных операций, ускоряет выпуск новых версий и делает процесс поставки программного продукта более прозрачным и контролируемым.

Для инженерных команд критически важны не только внутренние механизмы исполнения сборок, но и удобные средства мониторинга и оперативного управления: централизованное отображение текущих запусков, контроль сбоев, анализ логов, возможность ручного перезапуска упавших задач и оценка стабильности процессов. При использовании микросервисной архитектуры и распределении сервисов по десяткам репозиториев стандартные интерфейсы отдельных CI/CD-систем приводят к фрагментации контроля. Возникает потребность в специализированном веб-приложении, агрегирующем данные из внешних сервисов автоматизации в концепции единого окна («Single Pane of Glass»).

В рамках анализа предметной области рассмотрены наиболее распространенные решения: GitHub Actions, GitLab CI/CD, Jenkins и Argo CD.

\textbf{GitHub Actions} --- встроенная в платформу GitHub система автоматизации, в которой сценарии описываются в виде workflow-файлов формата YAML. Система глубоко интегрирована с событиями репозитория: pull requests, коммитами, релизами и внутренним хранилищем секретов. Данное решение оптимально для проектов, размещенных на GitHub~\cite{github_actions_docs}. Однако его ключевым ограничением является жесткая привязка к экосистеме GitHub, что затрудняет сквозной мониторинг гетерогенной инфраструктуры.

\textbf{GitLab CI/CD} --- компонент комплексной DevOps-платформы GitLab, конфигурируемый через файл \texttt{.gitlab-ci.yml}. Решение предоставляет развитые возможности по управлению стадиями, зависимостями задач, собственными runner-агентами, переменными окружения и артефактами~\cite{gitlab_cicd_docs}. Главным достоинством является полнота экосистемы, а ограничением --- ориентация на закрытый контур GitLab, требующая отдельной настройки внешних интеграций через REST API и webhooks.

\textbf{Jenkins} --- классическая расширяемая система автоматизации с открытым исходным кодом. Механизм Jenkins Pipeline позволяет описывать сложные сценарии сборки в файлах Jenkinsfile~\cite{jenkins_pipeline_docs}. Система независима от Git-хостинга и обладает обширной экосистемой плагинов. К ее недостаткам относятся высокая ресурсоемкость, сложность администрирования и визуальная неоднородность интерфейса.

\textbf{Argo CD} --- декларативный инструмент GitOps-доставки приложений в среду Kubernetes. Argo CD отслеживает состояние манифестов в Git-репозитории и синхронизирует их с фактическим состоянием кластера~\cite{argo_cd_docs}. Данный инструмент не предназначен для выполнения этапов компиляции и тестирования, а используется на финальной стадии доставки совместно с традиционными CI-системами.

Сравнительная характеристика рассмотренных инструментов автоматизации приведена в таблице~\ref{tab:cicd_compare}.

\begin{table}[H]
    \centering
    \caption{Сравнительный анализ CI/CD-систем}
    \label{tab:cicd_compare}
    \small
    \begin{tabular}{|p{2.4cm}|p{3.1cm}|p{3.1cm}|p{3.2cm}|p{3.2cm}|}
        \hline
        \textbf{Система} & \textbf{Основное назначение} & \textbf{Сильные стороны} & \textbf{Ограничения} & \textbf{Значимость для проекта} \\ \hline
        GitHub Actions & Автоматизация workflow внутри GitHub & Простая настройка, развитая экосистема actions, тесная связь с PR & Привязка к GitHub, лимиты бесплатного времени runner & Источник данных для интеграции через GitHub REST API \\ \hline
        GitLab CI/CD & Полный цикл CI/CD в среде GitLab & Единая платформа, runner-агенты, артефакты, environments & Зависимость от инфраструктуры GitLab & Источник данных для интеграции через GitLab API \\ \hline
        Jenkins & Универсальная система автоматизации & Предельная гибкость, независимость от хостинга кода & Сложность поддержки, устаревший базовый интерфейс & Потенциальный источник интеграции внешних jobs \\ \hline
        Argo CD & GitOps-доставка в Kubernetes & Синхронизация состояний, надежный деплой в кластер & Не выполняет этапы сборки и модульного тестирования & Пример специализированного инструмента доставки \\ \hline
    \end{tabular}
\end{table}

Проведенный анализ показывает, что инструменты CI/CD эффективно функционируют внутри своих экосистем, однако в командах с гетерогенной инфраструктурой инженерам приходится постоянно переключаться между различными панелями управления. 

Разрабатываемое веб-приложение призвано решить эту проблему в роли централизованного агрегатора и пульта управления: оно должно собирать актуальные статусы сборок через API внешних систем, хранить историю запусков, предоставлять быстрый доступ к логам и поддерживать диспетчеризацию (повторный запуск упавших пайплайнов) из единого интерфейса.

\subsection{Анализ программных средств разработки и выбор технологического стека}

Для реализации веб-приложения мониторинга и диспетчеризации пайплайнов необходим технологический стек, обеспечивающий высокую отзывчивость интерфейса, эффективную асинхронную обработку внешних сетевых запросов и безопасное хранение учетных записей и токенов интеграций.

Архитектура системы строится на разделении клиентской части (Single Page Application, SPA) и серверной части (REST API), взаимодействующих по протоколам HTTP/JSON и WebSocket.

\textbf{Клиентская часть.} Интерфейс системы мониторинга должен в реальном времени отображать статусы большого количества параллельных процессов, предоставлять средства динамической фильтрации, поиска по веткам и коммитам, а также отрисовывать графики стабильности сборок. Для разработки клиентского приложения рассмотрены технологии React, Vue.js и серверная генерация HTML (Jinja/Blade).

\begin{table}[H]
    \centering
    \caption{Сравнение фронтенд-технологий}
    \label{tab:frontend_compare}
    \small
    \begin{tabular}{|p{3cm}|p{3.7cm}|p{3.7cm}|p{3.7cm}|}
        \hline
        \textbf{Критерий} & \textbf{React (выбран)} & \textbf{Vue.js} & \textbf{Jinja (серверный рендеринг)} \\ \hline
        Архитектурный подход & SPA на базе компонентов и Virtual DOM~\cite{react_docs} & SPA на базе реактивных компонентов & Серверная генерация статических страниц \\ \hline
        Динамичность интерфейса & Мгновенное обновление карточек статусов без перезагрузки & Быстрое реактивное обновление DOM-дерева & Требует полной перезагрузки страницы или сложных скриптов \\ \hline
        Экосистема для дашбордов & Огромный выбор библиотек (Lucide, Recharts, Tailwind UI) & Достаточный выбор библиотек компонентов & Ограничена базовыми CSS-фреймворками \\ \hline
        Разделение ответственности & Полная изоляция логики представления от бэкенда & Полная изоляция клиентской части & Жесткая связность с кодом серверного фреймворка \\ \hline
    \end{tabular}
\end{table}

В качестве основы фронтенда выбрана библиотека \textbf{React}. Компонентный подход позволяет изолировать логику отображения отдельных карточек пайплайнов, индикаторов статусов и модальных окон управления. Использование Virtual DOM и хуков состояния обеспечивает высокую производительность при частых обновлениях данных через WebSockets без необходимости перерисовки всего экрана. Сборка проекта осуществляется с использованием современного инструмента Vite.

\textbf{Серверная часть.} Сервер приложения выполняет функции шлюза и агрегатора: принимает HTTP-запросы от клиентской части, осуществляет фоновый опрос и прием webhooks от внешних платформ (GitHub Actions, GitLab CI/CD), сохраняет промежуточные снапшоты в БД и транслирует команды повторного запуска (re-run). Для серверной реализации рассмотрены фреймворки на языке Python: FastAPI, Flask и Django.

\begin{itemize}
    \item \textbf{FastAPI (выбран)}: современный высокопроизводительный асинхронный фреймворк, построенный на базе ASGI-сервера Starlette и библиотеки валидации данных Pydantic~\cite{fastapi_docs}. Ключевым преимуществом FastAPI является нативная поддержка корутин (\texttt{async/await}), что критически важно для приложения-агрегатора, совершающего множество параллельных сетевых I/O-запросов к внешним API без блокировки основного потока. Кроме того, фреймворк автоматически генерирует интерактивную OpenAPI-документацию (Swagger UI).
    \item \textbf{Flask}: классический микрофреймворк с синхронной архитектурой (WSGI). Реализация высоконагруженного асинхронного взаимодействия с внешними API и работа с протоколом WebSocket во Flask требуют подключения дополнительных надстроек, что усложняет архитектуру.
    \item \textbf{Django}: полнофункциональный монолитный фреймворк. Включает избыточный для компактного API-сервиса функционал (административную панель, встроенный ORM с привязкой к синхронным вызовам), что создает лишние накладные расходы при разработке легковесного шлюза.
\end{itemize}

С учетом необходимости асинхронного ввода-вывода и строгой типизации контрактов API в качестве серверного фреймворка выбран \textbf{FastAPI}.

\textbf{Хранение данных.} Для долговременного хранения сведений о зарегистрированных пользователях, проектах, конфигурациях токенов доступа и исторических метрик сборок выбрана реляционная СУБД \textbf{PostgreSQL}~\cite{postgresql_docs}. Наличие поддержки типа данных \texttt{JSONB} позволяет сохранять неструктурированные метаданные внешних webhooks без усложнения реляционной схемы. Для взаимодействия с базой данных используется асинхронная ORM SQLAlchemy~2.0 в связке с драйвером \texttt{asyncpg}, а миграции схемы выполняются с помощью инструмента Alembic.

\textbf{Взаимодействие компонентов и внешняя интеграция.} Связь между SPA-клиентом и бэкендом организуется через REST API в формате JSON, а для моментальной доставки сведений об изменении статуса пайплайна используются протокол WebSocket и механизм Webhooks. Сетевые запросы к API внешних провайдеров (GitHub REST API v3, GitLab REST API v4) осуществляются с использованием асинхронного HTTP-клиента \texttt{httpx}.

Сводная информация о выбранном стеке приведена в таблице~\ref{tab:stack}.

\begin{table}[H]
    \centering
    \caption{Выбранный технологический стек}
    \label{tab:stack}
    \small
    \begin{tabular}{|p{4cm}|p{4.6cm}|p{5.2cm}|}
        \hline
        \textbf{Компонент} & \textbf{Выбранная технология} & \textbf{Назначение} \\ \hline
        Клиентская часть & React, TypeScript, Vite & Интерактивный SPA-интерфейс дашборда, карточки пайплайнов, графики \\ \hline
        Серверная часть & Python, FastAPI & Асинхронный REST API, диспетчеризация задач, валидация данных через Pydantic \\ \hline
        База данных & PostgreSQL & Надежное хранение проектов, пользователей, токенов и истории запусков \\ \hline
        ORM и миграции & SQLAlchemy 2.0 (async), Alembic & Управление реляционной схемой и асинхронный доступ к БД \\ \hline
        Обмен данными & REST API, Webhooks, WebSocket & Взаимодействие с клиентом и интеграция с внешними CI/CD-платформами \\ \hline
        Развертывание & Docker, Docker Compose & Контейнеризация сервисов и изоляция тестового окружения \\ \hline
    \end{tabular}
\end{table}

Выбранная архитектура на базе FastAPI и React разделяет ответственность между представлением и бизнес-логикой, гарантирует высокую скорость отклика и позволяет гибко подключать новых CI/CD-провайдеров.

\subsection{Формулировка цели, задач и требований проекта}

Основная инженерная проблема заключается в фрагментации данных о сборках при одновременной разработке множества микросервисов в разных системах контроля версий. Разрабатываемое веб-приложение призвано стать единой точкой наблюдения и диспетчеризации.

\textbf{Цель работы} --- разработать веб-приложение для централизованного управления и мониторинга CI/CD-пайплайнов, агрегирующее статусы сборок из внешних систем автоматизации и предоставляющее пользователю интерфейс оперативного реагирования и аналитики.

Для достижения поставленной цели необходимо решить следующие \textbf{задачи}:
\begin{enumerate}
    \item Провести анализ существующих платформ CI/CD и выявить ключевые функциональные требования к централизованной панели мониторинга.
    \item Обосновать выбор технологического стека для реализации SPA-клиента, асинхронного серверного API и хранилища данных.
    \item Спроектировать архитектуру веб-приложения и реляционную модель базы данных с учетом поддержки гетерогенных внешних источников.
    \item Реализовать модуль аутентификации пользователей и безопасного хранения токенов доступа к внешним сервисам.
    \item Разработать программные адаптеры для интеграции с внешними API (GitHub Actions REST API, GitLab API).
    \item Реализовать пользовательский интерфейс дашборда на React, обеспечивающий фильтрацию сборок, отображение карточек состояния и запуск команд повторного прогона (re-run).
    \item Реализовать механизмы оперативного обновления информации (Webhooks и WebSockets).
    \item Выполнить функциональное тестирование системы и оценить корректность обработки сбоев внешних сервисов.
\end{enumerate}

К разрабатываемому приложению предъявляются следующие \textbf{функциональные требования}:
\begin{itemize}
    \item аутентификация и авторизация пользователей в системе;
    \item управление подключениями к внешним провайдерам (ввод и безопасное хранение персональных токенов доступа);
    \item добавление и группировка отслеживаемых репозиториев и конвейеров сборки;
    \item сбор и синхронизация статусов пайплайнов (Success, Failed, Running, Queued, Canceled);
    \item отображение детальной информации о прогоне (автор коммита, ветка, хеш коммита, время выполнения и шаги пайплайна);
    \item возможность инициирования повторного запуска пайплайна (Re-run / Retry) через внешний API непосредственно из веб-интерфейса;
    \item фильтрация и поиск пайплайнов по статусу, проекту, автору и ветке;
    \item расчет базовых показателей надежности процессов (процент успешных прогонов и динамика длительности сборки).
\end{itemize}

Также сформированы \textbf{нефункциональные требования}:
\begin{itemize}
    \item время отдачи кэшированных агрегированных данных клиентскому приложению --- не более 300~мс;
    \item асинхронная неблокирующая архитектура взаимодействия с внешними сервисами;
    \item шифрование конфиденциальных данных (токенов доступа) при сохранении в базе данных;
    \item документирование REST API посредством автоматически генерируемой спецификации OpenAPI;
    \item масштабируемость архитектуры: возможность добавления новых CI/CD-провайдеров путем реализации единого программного интерфейса адаптера;
    \item адаптивный дизайн клиентской части для комфортной работы на экранах с различным разрешением;
    \item поддержка контейнеризации всех сервисов с помощью Docker Compose для быстрого развертывания.
\end{itemize}

Сформулированные требования определяют вектор дальнейшего проектирования системы: разработку диаграмм вариантов использования, проектирование базы данных и разработку макетов интерфейса в следующих разделах работы.